\documentclass[../main.tex]{subfiles}

\begin{document}

\chapter{Appendix A}

Here are the solutions to the practice problems at the end of each
notes.

\section{Solutions for Note 1}
\begin{enumerate}
\item
It suffices to show that $\alpha (x, y, z) + \beta (x', y', z') \in S$
for any scalars $\alpha$ and $\beta$. With standard vector addition
and scalar multiplication, $\alpha (x, y, z) + \beta (x', y', z') =
(0, \alpha y + \beta y', \alpha z + \beta z') \in S$. Thus, the set
is a vector space.

\item
Let $f$ and $g$ be functions in $\mathbb{R}$ such that $f'(0) = 0$
and $g'(0) = 0$. Notice that $(f + g)'(0) = f'(0) + g'(0) = 0$ and
$f + g \in S$. Moreover, for any scalars $\alpha$, $(\alpha f)'(0)
= \alpha f'(0) = 0$. Thus $\alpha f \in S$ and $S$ is a vector
space.

\item
First, for any scalars $\alpha$, notice that $\alpha(x_1, \ldots, x_n)
= (\alpha x_1, \ldots, \lambda x_n)$. If each coefficient is divided
by $\alpha \neq 0$, the scalar multiplication holds. The case where
$\alpha = 0$ naturally holds. Continuing, notice that $(x_1, \ldots,
x_n) + (x_1', \ldots, x_n') = (x_1 + x_1', \ldots, x_n + x_n') \in S$
due to the property of addition. Thus, $S$ is a vector space.

\item
Let $V = W_1 \cap \cdots \cap W_n$ where $W_i$ for integer $i \in
[1, n]$ is a subspace of $V$. For $u, v \in V$, notice that since
$u, v \in W_i$ for all $i$ and $\alpha u + \beta v \in W_i$ for any
scalars $\alpha$ and $\beta$ by construction, $\alpha u + \beta v \in
V$ and the intersection of subspaces of a vector space is also a
subspace of the vector space.

\item
First and foremost, it is evident that $e^{a_1 x}$ is linearly
independent since $e^{a_1 x} > 0$ for all $a_1$ and $x$. Therefore,
it suffices to show that if the set is linearly independent for
$n - 1$, then it is linearly independent for $n$. Consider the
following equations for some $c_k$.
\begin{align*}
    \sum_{k=1}^n c_k e^{a_k x} &= 0 \\
    \sum_{k=1}^n c_k a_k e^{a_k x} &= 0 \\
    \sum_{k=1}^n c_k a_n e^{a_k x} &= 0
\end{align*}
Subtracting the third equation, which is obtained by multiplying
$a_n$ in both sides, to the second equation obtained by
differentiating both sides,
\[
\sum_{k=1}^n c_k (a_k - a_n) e^{a_k x}
= \sum_{k=1}^{n-1} c_k (a_k - a_n) e^{a_k x}
= 0
\]
is obtained. Because $a_k \neq a_n$ for integer $k \in [1, n-1]$ and
$\{ e^{a_1 x}, \ldots, \linebreak e^{a_{n-1} x} \}$ is linearly
independent $c_1 = \cdots = c_{n-1} = 0$. In other words,
$c_n e^{a_n x} = 0$ and $c_n = 0$. Therefore,  $\{ e^{a_1 x}, \ldots,
e^{a_n x} \}$ is linearly independent by induction.

\item
By definition, note that $C = A + B$. Therefore, it suffices to show
that $A \cap B = \{ 0 \}$. Notice that $A$ and $B$ can be written as
$A = \Span\{ (1, 1, 1) \}$ and $B = \Span\{ (20, 15, 12) \}$.
Because $\lambda (1, 1, 1) = (20, 15, 12)$ if and only if $\lambda =
0$, $A \cap B = \{ 0 \}$ and $C = A \oplus B$.

\item
Consider the following equations for some scalars $k_1, k_2, k_3$.
\begin{align*}
    k_1 (a + b) + k_2 (b + c) + k_3 (c + a) &= 0 \\
    a (k_3 + k_1) + b (k_1 + k_2) + c (k_2 + k_3) &= 0
\end{align*}
The linear combinations of $\{ a + b, b + c, c + a \}$ equals zero
if and only if the linear combinations of $a$, $b$, and $c$ equals
zero. Moreover, from the equation above, $k_1 = -k_2 = k_3 = -k_1$.
In other words, $k_1 = k_2 = k_3 = 0$ and $\{ a + b, b + c, c + a \}$
is linearly independent.

\item
Let $f$ and $g$ be polynomials such that $f(0) = g(0) = 0$. Therefore,
$(f + g)(0) = 0$ and $f + g$ is in the set. Moreover, $(\alpha f)(0)
= 0$ for any scalars $\alpha$. Therefore, the set of all polynomials
in $\mathbb{P}^3$ that includes the origin is a subspace of
$\mathbb{P}^3$.
\end{enumerate}

\section{Solutions for Note 2}
\begin{enumerate}
\item
\begin{enumerate}
\item
Consider the following computation.
\begin{align*}
    A - 2B
    &=
    \begin{bmatrix}
        1 & 2 \\
        2 & 0
    \end{bmatrix}
    -
    2\begin{bmatrix}
        -1 & 2 \\
        0  & 3
    \end{bmatrix}
    =
    \begin{bmatrix}
        1 & 2 \\
        2 & 0
    \end{bmatrix}
    +
    \begin{bmatrix}
        2 & -4 \\
        0 & -6
    \end{bmatrix} \\
    &=
    \begin{bmatrix}
        3 & -2 \\
        2 & -6
    \end{bmatrix}
\end{align*}
\item
Consider the following equation.
\[
AB =
\begin{bmatrix}
    1 & 2 \\
    2 & 0
\end{bmatrix}
\begin{bmatrix}
    -1 & 2 \\
    0  & 3
\end{bmatrix}
=
\begin{bmatrix}
    -1 & 8 \\
    -2 & 7
\end{bmatrix}
\]
\item
Because $AB$ is found above, $BA$ could be found.
\[
BA =
\begin{bmatrix}
    -1 & 2 \\
    0  & 3
\end{bmatrix}
\begin{bmatrix}
    1 & 2 \\
    2 & 0
\end{bmatrix}
=
\begin{bmatrix}
    3 & -2 \\
    6 & 0
\end{bmatrix}
\]
Therefore, the following equation is true.
\begin{align*}
    (AB - BA)^\intercal
    &=
    \left(
    \begin{bmatrix}
        -1 & 8 \\
        -2 & 7
    \end{bmatrix}
    -
    \begin{bmatrix}
        3 & -2 \\
        6 & 0
    \end{bmatrix}
    \right)^\intercal
    =
    \left(
    \begin{bmatrix}
        -4 & 10 \\
        -8 & 7
    \end{bmatrix}
    \right)^\intercal \\
    &=
    \begin{bmatrix}
        -4 & -8 \\
        10 & 7
    \end{bmatrix}
\end{align*}
\end{enumerate}

\item
By the definition of matrix multiplication, $a_2B = [4, 5, 6]$.

\item
Notice that $A$ is a $2 \times 2$ square matrix. Therefore, define
$A$ as the following.
\[
A =
\begin{bmatrix}
    a_{11} & a_{12} \\
    a_{21} & a_{22}
\end{bmatrix}
\]
Therefore, the following equation is obtained.
\begin{align*}
2A - 3A^\intercal
&=
\begin{bmatrix}
    2a_{11} & 2a_{12} \\
    2a_{21} & 2a_{22}
\end{bmatrix}
-
\begin{bmatrix}
    3a_{11} & 3a_{21} \\
    3a_{12} & 3a_{22}
\end{bmatrix} \\
&=
\begin{bmatrix}
    2a_{11} - 3a_{11} & 2a_{12} - 3a_{21} \\
    2a_{21} - 3a_{12} & 2a_{22} - 3a_{22}
\end{bmatrix}
=
\begin{bmatrix}
    -1 & -5 \\
    0  & -4
\end{bmatrix}
\end{align*}
Continuing, the elements of $A$ could be found.
\begin{align*}
    2a_{11} - 3a_{11} &= -1 \\
    2a_{12} - 3a_{21} &= -5 \\
    2a_{21} - 3a_{12} &= 0  \\
    2a_{22} - 3a_{22} &= -4
\end{align*}
Therefore, $a_{11} = 1$, $a_{12} = 2$, $a_{21} = 3$, and $a_{22} = 4$.
Thus, $A$ can be defined as the following.
\[
A =
\begin{bmatrix}
    1 & 2 \\
    3 & 4
\end{bmatrix}
\]

\item
For any square matrices $S$, define $A$ and $B$ as the following.
\[
A = \frac{S + S^\intercal}{2}, \quad
B = \frac{S - S^\intercal}{2}
\]
Notice that $S = A + B$. Therefore, it suffices to show that $A$ is
symmetric and $B$ is skew-symmetric. By Theorem \ref{theo:Properties
of Transpose of a Matrix}, the following equations hold.
\[
A^\intercal
= \frac{S + S^\intercal}{2}, \quad
B^\intercal
= \frac{S^\intercal - S}{2}
= -B
\]
Thus, all square matrices can be represented as the sum of a symmetric
and skew-symmetric matrices.

\item
By definition, there exists $n, m \in \mathbb{Z}$ such that
$A^n = 0$ and $B^m = 0$. Because $AB = BA$, $A^{nm} B^{nm} = 0$ can
be written as $(AB)^{nm} = 0$ and $AB$ is nilpotent.

\item
By Theorem \ref{theo:Properties of Transpose of a Matrix}, $(AB +
(AB)^\intercal)^\intercal = (AB)^\intercal + AB = AB +
(AB)^\intercal$. Therefore, $AB + (AB)^\intercal$ is symmetric.

\item
An easy trap for this problem is that you might be tempted to
do $\det(A^2 - B^2) = \det(A + B) \det(A - B)$. However, notice that
this does not hold in general as $AB \neq BA$ in general.

Consider the following matrix $P$ and its inverse $P^{-1}$.
\[
P =
\begin{bmatrix}
    I & I \\
    I & -I
\end{bmatrix}, \quad
P^{-1} =
\frac{1}{2}
\begin{bmatrix}
    I & I \\
    I & -I
\end{bmatrix}
\]
Continuing,
\begin{align*}
    P^{-1}
    \begin{bmatrix}
        A & B \\
        B & A
    \end{bmatrix}
    P
    &= \frac{1}{2}
    \begin{bmatrix}
        A + B & B + A \\
        A - B & B - A
    \end{bmatrix}
    \begin{bmatrix}
        I & I \\
        I & -I
    \end{bmatrix} \\
    &=
    \begin{bmatrix}
        A + B & 0 \\
        0     & A - B
    \end{bmatrix}
\end{align*}
holds and
$\begin{bmatrix}
    A & B \\
    B & A
\end{bmatrix}
\cong
\begin{bmatrix}
    A + B & 0 \\
    0     & A - B
\end{bmatrix}$.
Therefore, the determinant of the original matrix is
$\det(A + B) \det(A - B)$.

\item
By definition, $A$ is skew-symmetric if $A^\intercal = -A$.
Notice that $(A^3)^\intercal = (A^\intercal)^3$ by Theorem
\ref{theo:Properties of Transpose of a Matrix}. Therefore,
$(A^3)^\intercal = (A^\intercal)^3 \linebreak = -A^3$. Thus, if $A$ is
skew-symmetric, then $A^3$ is also skew-symmetric.
\end{enumerate}

\section{Solutions for Note 3}
\begin{enumerate}
\item

\item

\item

\item

\item

\item

\end{enumerate}

\section{Solutions for Note 4}
\begin{enumerate}
\item

\item

\item

\item

\item

\item

\end{enumerate}
\end{document}


